% !TeX root = ../Book1.tex
% 一维叶片的二维示意；一般情形由r维切片替代曲线。
\begin{figure}[htbp]
	\centering
	\begin{tikzpicture}[font=\small]
		\begin{scope}
			\node at (1.8,3.45) {参数盒 \(B\times T\)};
			\fill[GeometryBlueLight] (0,0) rectangle (3.6,3);
			\draw[GeometryGray,dashed] (0,0) rectangle (3.6,3);
			\draw[orange!75!black,thick] (1.8,0.25)--(1.8,2.75);
			\foreach \h in {0.65,2.35}{
				\draw[GeometryGray] (0.25,\h)--(3.35,\h);
			}
			\draw[GeometryBlue,very thick,->] (0.25,1.5)--(3.35,1.5);
			\fill[GeometryBlue] (1.8,1.5) circle[radius=2pt];
			\node[above,fill=GeometryBlueLight,inner sep=1.5pt] at (2.75,1.5) {\(v=c\)};
			\node[orange!75!black,above] at (1.8,2.8) {\(u=0\)};
			\draw[->,GeometryGray] (0.25,-0.35)--(3.35,-0.35) node[right] {\(u\)};
			\draw[->,GeometryGray] (-0.35,0.25)--(-0.35,2.8) node[above] {\(v\)};
		\end{scope}
		\draw[->,thick] (4.05,1.5)--(5.35,1.5) node[midway,above] {\(\Psi\)};
		\begin{scope}[xshift=5.8cm]
			\node at (1.8,3.45) {流形中的坐标邻域 \(U\)};
			\fill[GeometryBlueLight] (0,0) rectangle (3.6,3);
			\draw[GeometryGray,dashed] (0,0) rectangle (3.6,3);
			\draw[orange!75!black,thick] (1.8,0.25)--(1.8,2.75);
			\foreach \h in {0.65,2.35}{
				\draw[GeometryGray] (0.25,\h) .. controls (1.1,{\h+0.45}) and (2.5,{\h-0.45}) .. (3.35,\h);
			}
			\draw[GeometryBlue,very thick] (0.25,1.5) .. controls (1.1,1.95) and (2.5,1.05) .. (3.35,1.5);
			\draw[GeometryBlue,thick,->] (2.05,1.44)--(2.7,1.37);
			\fill[GeometryBlue] (1.8,1.5) circle[radius=2pt];
			\node[GeometryBlue,above] at (2.55,1.47) {\(D\)};
			\node[GeometryBlue,right] at (3.4,1.5) {\(P_c\)};
			\node[orange!75!black,above] at (1.8,2.8) {\(\sigma(y)\)};
		\end{scope}
	\end{tikzpicture}
	\caption{一维叶片的二维示意。固定横向坐标的切片经 \(\Psi\) 变成积分叶片 \(P_c\)，橙色线表示横截片，蓝色箭头表示分布方向。高秩情形用 \(r\) 维片替代曲线，横向坐标仍沿片保持不变。}
	\label{g1:fig:frobenius-plaques}
\end{figure}
